% !TEX root = ../main.tex
\section[\texorpdfstring{\protect\fbox{\strut Quick Guide: Running the Notebook}\protect\hspace{6pt}\protect\colorbox{black}{\protect\color{white}\protect\bfseries\ START HERE\ }}{Quick Guide: Running the Notebook}]{Quick Guide: Running the Notebook}
\label{app:run}

This guide runs \code{CS24045.ipynb} from the submission \code{.zip} on a clean machine. It needs
Python 3.11 or newer and an internet connection for \code{pip install}.

\subsection{Files in the submission .zip}
The \code{.zip} holds exactly these nine files in one flat folder. The notebook looks for its input
files in its own folder (then \code{./data}, \code{../data}), so keep all files together and start
Jupyter from that folder.

\begingroup\small
\begin{tabularx}{\linewidth}{@{}lX@{}}
\toprule
\textbf{File} & \textbf{Purpose} \\
\midrule
\code{CS24045.ipynb} & the notebook (code, explanations, all four test runs) \\
\code{CS24045.pdf} & this report \\
\code{Environment.txt}, \code{Schedule.txt} & input files of test run 1 (our own) \\
\code{Environment\_Sample.txt}, \code{Schedule\_Sample.txt} & input files of test run 2 (provided sample) \\
\code{Environment\_Error.txt}, \code{Schedule\_Error.txt} & input files of test run 3 (failure case) \\
\code{requirements.txt} & dependencies needed to run the notebook (jupyter, matplotlib, tabulate) \\
\bottomrule
\end{tabularx}
\endgroup

\subsection{Windows 10/11 (PowerShell)}
\begin{enumerate}[itemsep=2pt]
\item Install Python 3.11 or newer from \code{python.org}. Tick ``Add python.exe to PATH''.
\item Unzip the submission and open PowerShell in that folder.
\item Create and activate a virtual environment:
\begin{Verbatim}[frame=single,fontsize=\small]
py -3 -m venv .venv
.venv\Scripts\Activate.ps1
\end{Verbatim}
If PowerShell blocks the activation script, run this first (it only affects this window):
\begin{Verbatim}[frame=single,fontsize=\small]
Set-ExecutionPolicy -Scope Process -ExecutionPolicy Bypass
\end{Verbatim}
\item Install the dependencies and start Jupyter:
\begin{Verbatim}[frame=single,fontsize=\small]
pip install -r requirements.txt
jupyter notebook CS24045.ipynb
\end{Verbatim}
\item In the browser, choose \textbf{Kernel $\rightarrow$ Restart \& Run All}.
\end{enumerate}
VS Code alternative: open the folder, select the \code{.venv} interpreter as the kernel, and click
\textbf{Run All}.

\newpage
\subsection{Linux (bash or zsh)}
\begin{enumerate}[itemsep=2pt]
\item Unzip the submission and open a terminal in that folder.
\item Create and activate a virtual environment (fish: \code{source .venv/bin/activate.fish}):
\begin{Verbatim}[frame=single,fontsize=\small]
python3 -m venv .venv
source .venv/bin/activate
\end{Verbatim}
\item Install the dependencies and start Jupyter:
\begin{Verbatim}[frame=single,fontsize=\small]
pip install -r requirements.txt
jupyter notebook CS24045.ipynb
\end{Verbatim}
\item In the browser, choose \textbf{Kernel $\rightarrow$ Restart \& Run All}.
\end{enumerate}
On Debian or Ubuntu, if \code{venv} is missing: \code{sudo apt install python3-venv}.

\subsection{What you should see}
The notebook runs four test runs. Runs 1 and 2 end with \code{Status: SUCCESS} (total cost
64 and 34). Run 3 ends with \code{FAILURE: no path from D4 (9, 5) to D5 (9, 0)}.
This failure is intended: it is the failure-case test (Section~6).
Run 4(3b) ends with \code{FAILURE: CSP solver found no feasible schedule}. This failure is also intended.

\subsection{Source repository}
The full project is public at \code{https://github.com/Silicon-Triad/grid-sweeper}.
\begin{Verbatim}[frame=single,fontsize=\small]
git clone https://github.com/Silicon-Triad/grid-sweeper.git
cd grid-sweeper
pip install -r requirements.txt
pytest tests/
python tools/build_notebook.py --check
\end{Verbatim}
Where to look: \code{src/cep/} (the code), \code{tests/} (the test suite), \code{data/} (input
files), \code{docs/decisions.md} (design decisions \dec{01}, \dec{02}, \dots).
